La definición que hemos visto de objeto libre (\ref{def:free-object}) está
asociada a constructos y por tanto a la categoría $\bSet$. Sin embargo, no hay
razón para limitarse a este caso, y de hecho los objetos libres son un caso
particular de flecha universal en el que el codominio del funtor olvidadizo es
$\bSet$. Cuando todos los elementos de dicho codominio admiten una flecha
universal podemos definir un funtor libre asociado al funtor olvidadizo, y este
funtor tiene propiedades interesantes que conviene estudiar. Este capítulo se
basa principalmente en \cite[II.5, III.1--2 y IV.1]{maclane}.%TODO Citar todo

\section{Flechas universales}

\begin{definition}
  Sean $U:\cC\to\cB$ un funtor y $b$ un objeto de $\cB$, una \conc{flecha
    universal} de $b$ a $U$ es un par $(c,u)$ formado por un objeto $c$ de
  $\cC$, el \conc{objeto libre}, y un morfismo $u:b\to Uc$ en $\cB$, tales que
  para todo morfismo $f:b\to Ux$ en $\cB$ existe un único morfismo
  $\hat f:c\to x$ en $\cC$ tal que $f=U\hat f\circ u$, como se muestra en la figura
  \ref{fig:universal}.
  \begin{figure}
    \centering
    \begin{diagram}
      \path (0,2) node(B){$b$} (2,2) node(UC){$Uc$} (4,2) node(C){$c$};
      \path                    (2,0) node(UX){$Ux$} (4,0) node(X){$x$};
      \draw[->] (B) -- node[above]{$u$} (UC);
      \draw[->] (B) -- node[left]{$f$} (UX);
      \draw[->,dotted] (UC) -- node[right]{$U\hat f$} (UX);
      \draw[->,dotted] (C) -- node[right]{$\hat f$} (X);
    \end{diagram}
    \caption{Flecha universal de un objeto a un funtor.}
    \label{fig:universal}
  \end{figure}
\end{definition}

\begin{example}
  Las flechas universales suelen representar inmersiones de objetos en un
  cierto objeto completado o con estructura adicional.
  \begin{enumerate}
  \item Un objeto libre sobre un conjunto $X$ en un constructo es una flecha
    universal de $X$ al constructo.
  \item Si $U:\bField\inTo\bDom$ es el funtor inclusión de la subcategoría
    completa de los cuerpos en la categoría de dominios, una flecha universal de
    un dominio $D$ en $U$ es el conjunto cociente $Q(D)$ junto con la inclusión.
  \item Consideremos la categoría $\bMGrph$ de los \conc{multigrafos}, los
    grafos dirigidos (no necesariamente finitos) que admiten varios ejes entre
    dos mismos vértices. Si $U:\bCat\to\bMGrph$ es el funtor que <<olvida>> la
    composición y la identidad, un objeto libre de un multigrafo $M$ a $U$ es
    una categoría cuyos objetos son los vértices de $M$ y cuyos morfismos entre
    dos objetos son los caminos entre ellos en $M$, tomando como composición la
    concatenación de caminos y como identidad el camino vacío.
  \item Si $D:\cS\to\cC$ es un diagrama, un morfismo $f:D\to\Delta c$ en
    $\cC^\cS$ es un sumidero con codominio $c$ que conmuta con $D$, de modo que
    un colímite de $D$ es una flecha universal de $D$ a $\Delta:\cC\to\cC^\cS$.
  \end{enumerate}
\end{example}

Nos gustaría ver que las flechas universales son únicas salvo isomorfismo, pero
para ello primero tenemos que ver qué significa esto. Una forma de hacerlo es
definir la categoría en la que <<viven>> estas flechas.

\begin{definition}
  Si $U:\cC\to\cB$ es un funtor y $b$ es un objeto de $\cB$, la \conc{categoría
    de objetos $U$-bajo $b$}, $(b\downarrow U)$, tiene como objetos los pares
  $(c,f)$ formados por un objeto $c$ de $\cC$ y un morfismo $f:b\to Uc$; como
  morfismos $h:(c,f)\to(c',f')$ los morfismos $h:c\to c'$ en $\cC$ tales que
  $f'=Uh\circ f$, y como composición e identidad las correspondientes en $\cC$.
\end{definition}

\begin{proposition}
  La flecha universal de un objeto $b$ de $\cB$ a un funtor $U:\cC\to\cB$, si
  existe, es única salvo isomorfismo en $(b\downarrow U)$.
\end{proposition}
\begin{proof}
  Es el objeto inicial de $(b\downarrow U)$.
\end{proof}

El concepto dual al de flecha universal de un objeto a un funtor es el de flecha
universal de un funtor a un objeto.

\begin{definition}
  Sean $T:\cC\to\cB$ un funtor y $b$ un objeto de $\cB$, una \conc{flecha
    universal} de $T$ a $b$ es un par $(c,v)$ formado por un objeto $c$ de $\cC$
  y un morfismo $v:Tc\to b$ en $\cB$, tales que para todo morfismo $f:Tx\to b$
  en $\cB$ existe un único morfismo $\hat f:x\to c$ tal que $f=v\circ T\hat f$,
  como se muestra en la figura \ref{fig:couniversal}.
  
  \begin{figure}
    \centering
    \begin{diagram}
      \path (0,-2) node(B){$b$} (-2,-2) node(UC){$Tc$} (-4,-2) node(C){$c$};
      \path                     (-2,0) node(UX){$Tx$}  (-4,0) node(X){$x$};
      \draw[<-] (B) -- node[below]{$v$} (UC);
      \draw[<-] (B) -- node[right]{$f$} (UX);
      \draw[<-,dotted] (UC) -- node[left]{$T\hat f$} (UX);
      \draw[<-,dotted] (C) -- node[left]{$\hat f$} (X);
    \end{diagram}
    \caption{Flecha universal de un funtor a un objeto.}
    \label{fig:couniversal}
  \end{figure}
\end{definition}

\begin{example}
  Si $D:\cS\to\cC$ es un diagrama, un morfismo $f:\Delta c\to D$ en $\cC^\cS$ es
  una fuente con dominio $c$ que conmuta con $D$, con lo que un límite de $D$ es
  una flecha universal de $\Delta$ a $D$.
\end{example}

\begin{definition}
  Si $T:\cC\to\cB$ es un funtor y $b$ es un objeto de $\cB$, la \conc{categoría
    de objetos $T$-sobre $b$}, $(T\downarrow b)$, tiene como objetos los pares
  $(c,f)$ formados por un objeto $c$ de $\cC$ y un morfismo $f:Tc\to b$; como
  morfismos $h:(c,f)\to(c',f')$ los morfismos $h:c\to c'$ en $\cC$ tales que
  $f=f'\circ Th$.
\end{definition}

\begin{proposition}
  La flecha universal de un objeto $b$ de $\cB$ a un funtor $T:\cC\to\cB$, si
  existe, es única salvo isomorfismo en $(T\downarrow b)$.
\end{proposition}

\section{Lema de Yoneda}

El lema de Yoneda es un resultado clásico sobre transformaciones naturales que
permite relacionar las mismas con transformaciones naturales. Antes de verlo
conviene definir algunos funtores útiles.

\begin{definition}
  Si $\cC$ es una categoría con conjuntos hom pequeños, definimos el
  \conc{bifuntor hom} $\hom_\cC:\dual{\cC}\times\cC$ sobre objetos $(a,b)$ como
  $\hom_\cC(a,b)$, y sobre morfismos $(f,g):(a,b)\to(a',b')$ como
  $\hom_\cC(f,g)(h)\coloneqq g\circ h\circ f$. Dados dos funtores $S:\cA\to\cC$
  y $T:\cB\to\cC$, definimos el bifuntor
  $\hom_\cC(S-,T-):\dual{\cA}\times\cB\to\cC$ como $\hom_\cC\circ(S\times T)$.
\end{definition}

\begin{definition}
  Dadas dos categorías $\cC$ y $\cD$, definimos el \conc{funtor de evaluación}
  $E:\cD^\cC\times\cC\to\cD$ sobre objetos como $E(T,c)\coloneqq Tc$ y sobre
  morfismos $(\tau,f):(T,c)\to(U,c')$ como
  $E(\tau,f)\coloneqq \tau f\coloneqq\tau_{c'}\circ Tf=Uf\circ\tau_c$.
\end{definition}

\begin{definition}
  Si $\cC$ es una categoría con conjuntos hom pequeños, llamamos \conc{funtor de
    Yoneda} al funtor $Y:\dual{\cC}\to\bSet^\cC$ que lleva objetos $c$ a
  funtores $\hom(c,-)$ y morfismos $f:a\to b$ en $\cC$ a transformaciones
  naturales $\hom(f,-):\hom(b,-)\to\hom(a,-)$ dadas por
  $\hom(f,-)_c(g)\coloneqq g\circ f$.
\end{definition}

\begin{lemma}[Yoneda]\label{lem:yoneda}
  Sea $\cC$ una categoría con conjuntos hom pequeños. Para cada funtor $T:\cC\to\bSet$
  y objeto $c$ de $\cC$, la función
  \[
    \gamma_{T,c}:\hom_{\bSet^\cC}(\hom_\cC(c,-),T)\to Tc
  \]
  dada por $\gamma_{T,c}(\tau)\coloneqq\tau_c1_c$ es una biyección natural entre
  el funtor $\hom(Y-,-)$ (con el orden de las entradas cambiado) y
  el funtor de evaluación $\bSet^\cC\times\cC\to\bSet$.
\end{lemma}
\begin{proof}
  \begin{figure}
    \hfil
    \begin{subfigure}{0.45\textwidth}
      \centering
      \begin{diagram}
        \path (0,2) node(UL){$\hom_\cC(c,c)$} (3,2) node(UR){$Tc$};
        \path (0,0) node(BL){$\hom_\cC(c,x)$} (3,0) node(BR){$Tx$};
        \draw[->] (UL) -- node[above]{$\tau_c$} (UR);
        \draw[->] (BL) -- node[above]{$\tau_x$} (BR);
        \draw[->] (UL) -- node[left]{$\hom_\cC(c,f)$}  (BL);
        \draw[->] (UR) -- node[right]{$Tf$} (BR);
      \end{diagram}
    \end{subfigure}
    \hfil
    \begin{subfigure}{0.45\textwidth}
      \centering
      \begin{diagram}
        \path (0,2) node(UL){$1_c$} (3,2) node(UR){$\tau_c1_c$};
        \path (0,0) node(BL){$f$}   (3,0) node(BR){$\tau_xf=(Tf)(\tau_c1_c)$};
        \tikzsquig(UL)--(UR); \tikzsquig(BL)--(BR);
        \tikzsquig(UL)--(BL); \tikzsquig(UR)--(BR);
      \end{diagram}
    \end{subfigure}
    \hfil
    
    \caption{Representación del lema de Yoneda.}
    \label{fig:yoneda}
  \end{figure}

  Para la biyección basta observar la figura \ref{fig:yoneda}. Para la
  naturalidad, fijado un morfismo $(\varphi,f):(S,a)\to(T,b)$ de
  $\bSet^\cC\times\cC$, debemos comprobar que el siguiente diagrama conmuta.
  
  \begin{center}
    \begin{diagram}
      \path (0,2) node(UL){$\hom(\hom(a,-),S)$} (3,2) node(UR){$Sa$};
      \path (0,0) node(BL){$\hom(\hom(b,-),T)$} (3,0) node(BR){$Tb$};
      \draw[->] (UL) -- node[above]{$\gamma_{S,a}$} (UR);
      \draw[->] (BL) -- node[above]{$\gamma_{T,b}$} (BR);
      \draw[->] (UL) -- node[left]{$\hom(\hom(f,-),\varphi)$} (BL);
      \draw[->] (UR) -- node[right]{$\varphi f$} (BR);
    \end{diagram}
  \end{center}

  Sea entonces $\tau:\hom(a,-)\to S$ una transformación natural,
  \begin{multline*}
    (\varphi f)(\gamma_{S,a}(\tau))
    = (\varphi f)(\tau_a1_a)
    = \varphi_b((Sf)(\tau_a1_a))
    = \varphi_b(\tau_b\hom(a,f)(1_a))
    = \varphi_b\tau_bf = \\
    = \gamma_{T,b}(\varphi\cdot\tau\cdot\hom(f,-))
    = \gamma_{T,b}(\hom(\hom(f,-),\varphi)(\tau)).
  \end{multline*}
\end{proof}

Este lema permite demostrar la siguiente relación entre isomorfismos naturales y
flechas universales.

\begin{proposition}\label{prop:yoneda-prop}
  Dado un funtor $U:\cC\to\cB$ y un objeto $b$ de $\cB$, un par $(c,u:b\to Uc)$
  es una flecha universal de $b$ a $U$ si y sólo si, para cada objeto $x$ en
  $\cC$, la función $\tau_x:\hom_\cC(c,x)\to\hom_\cB(b,Ux)$ dada por
  $\tau_x(f)\coloneqq Uf\circ u$ es biyectiva, en cuyo caso $\tau$ es un
  isomorfismo natural. Además, dados objetos $b$ de $\cB$ y $c$ de $\cC$, todo
  isomorfismo natural \[
    \hom_\cC(c,-)\cong\hom_\cB(b,U-)
  \]
  es de esta forma para un único morfismo $u:b\to Uc$ para el que $(c,u)$ es una
  flecha universal.
\end{proposition}
\begin{proof}
  La definición de flecha universal equivale a esta biyección, que es natural ya
  que, para cada morfismo $g:x\to y$ en $\cC$ y $f:c\to x$,
  $\hom(b,Ug)(\tau_x(f))=Ug\circ Uf\circ u=U(g\circ f)\circ
  u=U(\hom(c,g)(f))=\tau_y(\hom(c,g)(f))$.

  Para el recíproco, si $\tau:\hom_\cC(c,-)\to\hom_\cB(b,U-)$ es un isomorfismo
  natural, por el lema de Yoneda y la biyectividad de $\tau_x$ se tiene que todo
  morfismo $b\to Ux$ se expresa de forma única como
  $\tau_xf=\hom(b,Uf)(\tau_c1_c)=Uf\circ\tau_c1_c$ para cierto $f:c\to x$, lo
  que significa precisamente que $\tau_c1_c$ es universal de $b$ a $U$.
\end{proof}

\section{Adjunciones}

Sea $U:\cC\to\cB$ un funtor tal que todo objeto $b$ en $\cB$ admite una flecha
universal $(Fb,u_b)$ de $b$ a $U$. Para cada morfismo $f:b\to b'$ en $\cB$,
siguiendo la figura \ref{fig:universal} existe un único morfismo $Ff:Fb\to Fb'$
en $\cC$ tal que $UFf\circ u_b=u_{b'}\circ f$, y además claramente $F1_b=1_{Cb}$
y $Fg\circ Ff=F(g\circ f)$, de modo que $F:\cB\to\cC$ es un funtor y
$u:1_\cB\to U\circ F$ es una transformación natural.

El que la imagen de $u$ esté formada por flechas universales permite obtener una
especie de inversa. Dados un objeto $b$ en $\cB$ y un objeto $c$ en $\cC$, para
cada morfismo $f:b\to Uc$ existe un único $\hat f:Fb\to c$ tal que
$f=U\hat f\circ u_b$, de modo que $\psi_{b,c}:\hom_\cC(Fb,c)\to\hom_\cB(b,Uc)$
dada por $\psi_{b,c}(g)\coloneqq Ug\circ u_b$ es una biyección, y como $U$ es un
funtor y $u$ es natural, $\psi$ es un isomorfismo natural entre los funtores
$\hom_\cC\circ(F\times 1),\hom_\cB\circ(1\times U):\dual{\cB}\times\cC\to\bSet$
(si los conjuntos hom no son siempre pequeños, basta sustituir $\bSet$ por una
clase de conjuntos más grande).  Entonces $u_b\equiv\psi_{b,Fb}(1)$ y, del mismo
modo, podemos definir la transformación natural $e:F\circ U\to 1_\cC$ como
$e_c\coloneqq\psi_{Uc,c}^{-1}(1)$, de modo que para cada objeto $c$, $(Uc,e_c)$
es una flecha universal de $F$ a $c$.

La relación entre las transformaciones naturales $u$ y $e$ es más estrecha que
esto. Para un objeto $c$, $1_{Uc}=\psi(e_c)=Ue_c\circ u_{Uc}$, y de forma dual,
para un objeto $b$,
$\psi(e_{Fb}\circ Fu_b)=Ue_{Fb}\circ Ue_{Fb}\circ UFu_b\circ u_b=Ue_{Fb}\circ
u_{UFb}\circ u_b=1_{UFb}\circ u_b=u_b$ y por tanto
$1_{Fb}=\psi^{-1}(u_b)=e_{Fb}\circ Fu_b$.

Estas dos identidades se pueden expresar más elegantemente con la notación
adecuada. Si $\tau:R\to S$ es una transformación natural entre dos funtores
$R,S:\cB\to\cC$ y $T:\cC\to\cD$ es otro funtor, podemos definir la
transformación natural $T\tau:T\circ R\to T\circ S$ como
$(T\tau)_b\coloneqq T(\tau_b)$ para cada objeto $b$ en $\cB$. Por otro lado, si
$U:\cA\to\cB$ es otro funtor, podemos definir la transformación natural
$\tau U:R\circ U\to G\circ U$ como $(\tau U)_a\coloneqq\tau_{Ka}$ para cada
objeto $a$ en $\cA$.

Con todo esto en mente, definimos las adjunciones como sigue.

\begin{definition}
  Una \conc{adjunción} entre dos categorías $\cB$ y $\cC$ es una tupla
  $(F,G,\eta,\eps)$ formada por dos funtores $F:\cB\to\cC$ y $G:\cC\to\cB$ y dos
  transformaciones naturales $\eta:1\to GF$ y $\eps:FG\to 1$, llamadas
  respectivamente \conc{unidad} y \conc{co-unidad}, tales que
  $G\eps\cdot\eta G=1_G$ y $\eps F\cdot F\eta=1_F$.
\end{definition}

Cabe preguntarse si todas las adjunciones se pueden construir como en el
razonamiento anterior. La respuesta es que sí, como vemos en el siguiente
teorema.

\begin{theorem}\label{thm:adjoint-elems}
  Una adjunción $(F,G,\eta,\eps)$ entre $\cB$ y $\cC$ viene determinada por
  cualquiera de las siguientes listas de elementos:
  \begin{enumerate}
  \item \label{enu:adj-canon} Funtores $F$ y $G$ y un isomorfismo natural
    $\psi:\hom\circ(F\times 1)\to\hom\circ(1\times G)$. Entonces se definen
    $\eta_b\coloneqq\psi_{b,Fb}(1)$ y $\eps_c\coloneqq\psi_{Uc,c}^{-1}(1)$.
  \item \label{enu:adj-univ} El funtor $G$ y, para cada objeto $b$ en $\cB$, una
    flecha universal $(c_b,\eta_b)$ de $b$ a $G$. Entonces $F$ se define sobre
    cada objeto $b$ como $c_b$ y sobre cada morfismo $f:b\to b'$ como el único
    $Ff$ tal que $GFf\circ\eta_b=\eta_{b'}\circ f$, y el isomorfismo
    $\psi_{b,c}$ del apartado \ref{enu:adj-canon} se define como
    $g\mapsto Gg\circ\eta_b$.
  \item \label{enu:adj-couniv} El funtor $F$ y, para cada objeto $c$ en $\cC$,
    una flecha universal $(b_c,\eps_c)$ de $F$ a $b$. Entonces $G$ se define
    sobre cada objeto $c$ como $b_c$ y sobre cada morfismo $g:c\to c'$ como el
    único $Gg$ tal que $\eps_c\circ FGg=g\circ\eps_{c'}$, y el isomorfismo
    $\psi_{b,c}$ se define como el inverso de $f\mapsto\eps_c\circ Ff$.
  \end{enumerate}
  En particular, cada $(Fb,\eta_b)$ es una flecha universal de $b$ a $G$ y cada
  $(Gc,\eps_c)$ es una flecha universal de $F$ a $u$.
\end{theorem}
\begin{proof}
  Para (\ref{enu:adj-canon}), si $b$ es un objeto de $\cB$ y $c$ uno de $\cC$,
  definimos $\psi:\hom(Fb,c)\to\hom(b,Gc)$ como $\psi(g)\coloneqq Gg\circ\eta_b$
  y $\theta:\hom(b,Gc)\to\hom(Fb,c)$ como $\theta(f)\coloneqq\eps_c\circ
  Ff$. Entonces, como $\eta$ es natural,
  \[
    \psi(\theta(f))=G\eps_c\circ GFf\circ\eta_b=G\eps_c\circ\eta_{Gc}\circ f=f,
  \]
  por lo que $\psi\circ\theta=1$, y análogamente $\theta\circ\psi=1$, por lo que
  $\psi$ es un isomorfismo, claramente natural respecto a $b$ y $c$, y se tiene
  $\psi(1)=\eta_b$ y $\theta(1)=\eps_c$. Para el recíproco, definiendo $\eta_b$
  y $\eps_c$ a partir de $\psi$ como en el enunciado, para todo morfismo
  $f:Fb\to c$ podemos seguir flechas como sigue.

  \hfil
  \begin{diagram}
    \path (0,2) node(UL){$\hom(Fb,Fb)$} (3,2) node(UR){$\hom(b,GFb)$};
    \path (0,0) node(BL){$\hom(Fb,c)$}  (3,0) node(BR){$\hom(b,Gc)$};
    \draw[->] (UL) -- node[above]{$\psi_{b,Fb}$} (UR);
    \draw[->] (BL) -- node[above]{$\psi_{b,c}$} (BR);
    \draw[->] (UL) -- node[left]{$\hom(Fb,f)$} (BL);
    \draw[->] (UR) -- node[right]{$\hom(b,Gf)$} (BR);
  \end{diagram}
  \hfil
  \begin{diagram}
    \path (0,2) node(UL){$1$} (3,2) node(UR){$\eta_b$};
    \path (0,0) node(BL){$f$} (3,0) node(BR){$\psi_{b,c}(f)=Gf\circ\eta_b$};
    \tikzsquig (UL) -- (UR); \tikzsquig (BL) -- (BR);
    \tikzsquig (UL) -- (BL); \tikzsquig (UR) -- (BR);
  \end{diagram}
  \hfil

  Con esto $1_{Gc}=\psi(\eps_c)=G\eps_c\circ\eta_{Gc}$, y la otra identidad es
  dual a esta y se demuestra de forma análoga.

  Para (\ref{enu:adj-univ}), $(Fb,\eta_b)$ es una flecha universal, pues para
  cada objeto $x$ en $\cC$, $\psi_{b,x}(g)=Gg\circ\eta_b$ es una biyección
  $\hom(Fb,x)\to\hom(b,Ux)$ natural respecto a $x$
  (\ref{prop:yoneda-prop}). Recíprocamente, si sólo tenemos $G$ y las flechas
  universales $(c_b,\eta_b)$, $F$ definido de esta forma es un funtor que hace a
  $\eta$ natural y $\psi_{b,c}$ definido de esta forma es un isomorfismo por
  (\ref{prop:yoneda-prop}) y es claramente natural.

  Finalmente, (\ref{enu:adj-couniv}) es dual a (\ref{enu:adj-univ}).
\end{proof}

\begin{example}
  Este teorema permite definir una gran variedad de adjunciones.
  \begin{enumerate}
  \item El caso más sencillo de adjunción se da cuando todos los conjuntos
    admiten un objeto libre en un cierto constructo $\cC$. Entonces tenemos una
    adjunción $(F,U,\eta,\eps)$, donde $U:\cC\to\bSet$ es el funtor olvidadizo,
    $F:\bSet\to\cC$ es el funtor libre, $\eta_X:X\inTo UFX$ es la inclusión de
    la base en el conjunto subyacente del objeto y $\eps_c:FUc\epicTo c$ es el
    epimorfismo que aparece al describir un objeto como cociente de un cierto
    objeto libre, que resulta ser el que tiene los propios elementos de $c$ como
    generadores. Por ejemplo, en el caso de $R\dash\bMod$, $\eps_M$ lleva sumas
    formales $\sum_{i=1}^ka_im_i$, con cada $a_i\in R$ y cada $m_i\in M$, a su
    evaluación en el módulo $M$.
  \item Entre $\bDom$ y $\bField$ hay una adjunción $(Q,U,\eta,\eps)$ formada
    por la creación de cuerpos de fracciones $Q:\bDom\to\bField$, la inclusión
    de categorías $U:\bField\inTo\bDom$, la inclusión canónica $\eta_D:D\to UQX$
    y la identidad $\eps_K:QUK\to K$ (recordemos que el cuerpo de fracciones de
    un cuerpo es el propio cuerpo).
  \item Si $\cC$ es una categoría que tiene colímites con diagrama $\cS$, entre
    $\cC^\cS$ y $\cC$ hay una adjunción
    $(\underrightarrow{\lim},\Delta,\eta,\eps)$, donde $\underrightarrow{\lim}$
    lleva cada diagrama a su objeto colímite, $\Delta:\cC\to\cC^\cS$ es el
    funtor diagonal y $\eta_D:D\to\Delta\underrightarrow{\lim}D$ es el sumidero
    colímite. Respecto a $\eps_c:\underrightarrow{\lim}\Delta c\to c$, el
    colímite de $\Delta c$ es $^Ic$, siendo $I$ el número (cardinal) de
    componentes conexas de $\cS$, y $\eps_c$ es el morfismo que <<une todas las
    copias de $c$>>.
  \item Análogamente, si $\cC$ tiene límites con diagrama $\cS$, tenemos una
    adjunción $(\Delta,\underleftarrow{\lim},\eta,\eps)$ donde
    $\underleftarrow{\lim}$ lleva cada diagrama a su límite, $\eta_c$ es el
    límite de $\Delta c$ (una potencia de $c$) y $\eps_D$ es el límite como
    fuente.
  \end{enumerate}
\end{example}

\section{Adjuntos laterales}

\begin{definition}
  Dada una adjunción $(F,G,\eta,\eps)$, decimos que $F$ es un \conc{adjunto
    izquierdo} de $G$ y que $G$ es un \conc{adjunto derecho} de $F$.
\end{definition}

\begin{example}
  El funtor olvidadizo $U:\bTop\to\bSet$ tiene una adjunción por cada
  lado:\cite[4.1]{riehl}
  \begin{enumerate}
  \item Por la izquierda tiene el funtor $D:\bSet\to\bTop$ que asocia a cada
    conjunto su topología discreta. Entonces $\eta_X:X\to UDX$ es la identidad
    en $X$ y $\eps_T:DUT\to T$ es la identidad hacia $T$ desde el mismo conjunto
    pero con la topología discreta.
  \item Por la derecha tiene el funtor $N:\bSet\to\bTop$ que asocia a cada
    conjunto su topología discreta. Aquí $\eps_X:UNX\to X$ es la identidad en
    $X$ y $\eta_T:T\to NUT$ es la identidad desde $T$ hacia el mismo conjunto
    pero con la topología indiscreta.
  \end{enumerate}
\end{example}

\begin{proposition}
  El adjunto por la izquierda o por la derecha de un funtor es único salvo
  isomorfismo natural.
\end{proposition}
\begin{proof}
  El adjunto por la izquierda de un funtor $G:\cC\to\cB$ viene dado por una
  flecha universal $(Fb,\eta_b)$ de $b$ a $G$ para cada objeto $b$ en $\cB$
  (\ref{thm:adjoint-elems}), pero estas flechas son únicas salvo
  isomorfismo. Así, si $(Fb,\eta_b)_b$ y $(F'b,\eta'_b)_b$ son familias de estas
  flechas, para cada $b$ existe un único isomorfismo $h_b:Fb\to F'b$ tal que
  $\eta'_b=Gh_b\circ\eta_b$. Para ver que este es natural, para $f:b\to b'$, por
  la construcción de $Ff$ en \ref{thm:adjoint-elems},
  \[
    G(h_{b'}\circ Ff\circ h_b^{-1})\circ\eta'_b = Gh_{b'}\circ GFf\circ\eta_b
    = Gh_{b'}\circ GFf\circ\eta_b = Gh_{b'}\circ\eta_{b'}\circ f = \eta'_{b'}\circ f,
  \]
  y por la unicidad en dicha construcción, $F'f=h_{b'}\circ Ff\circ h_b^{-1}$.

  El concepto de adjunto por la derecha es el dual.
\end{proof}

%%% Local Variables:
%%% mode: latex
%%% TeX-master: "main"
%%% End:
